\section{Microscopic stability above the mean-spacing error scale}\label{sec:robustness}
This addendum strengthens the consequences of \hyperref[sec:locality]{fixed-strength locality}. The finite covering bounds below are elementary. Their role is to make the perturbation and unfolding thresholds explicit; no priority claim is made. In particular, operator-norm accuracy better than the mean spacing is \textbf{not} necessary for the negative local-statistics conclusion. The unperturbed reference still has arithmetic clusters. Nothing here proves positive BGS for a nonarithmetic chaotic system.

\subsection{A finite covering theorem}\label{robustness:a-finite-covering-theorem}

Let \(x_1,\ldots,x_D\) be a multiset on \(\mathbb R/\mathbb Z\), including multiplicities, covered by at most \(K\) arcs whose total length is at most \(\ell\). Assume \(0\leq\ell<1\). No uniform placement or equal occupancy of the arcs is assumed. Write \(S_j\) for the cyclic gaps multiplied by \(D\), and set \[
 R_D(x)=\frac1D\#\{(i,j):i\ne j,\ D\operatorname{dist}(x_i,x_j)\leq x\}.
\] For every \(x>0\) and \(0<L<D\), \[
 \frac1D\#\{j:S_j>x\}\leq\frac KD+\frac{\ell}{x},          \tag{RB.R1}
\] \[
 R_D(x)\geq\frac1{K/D+\ell/x}-1,                           \tag{RB.R2}
\] and, for a uniformly distributed origin \(z\), \[
 \mathbb P\{[z,z+L/D)\text{ contains a point}\}
                   \leq KL/D+\ell.                        \tag{RB.R3}
\] The last upper bound can of course be replaced by its minimum with one. The denominator in (RB.R2) is positive for a nonempty multiset; \(K\geq1\).

\textbf{Proof.} Merge overlapping arcs first. There are at most \(K\) components, with total length at most \(\ell\), and a point outside their union can be chosen as the seam of the circle. There are at most \(K\) gaps connecting different components. The lengths of the remaining gaps sum to at most \(\ell\). Each gap counted in (RB.R1) has length greater than \(x/D\), proving (RB.R1), without any lower bound on individual gaps.

Divide each covering arc into consecutive cells of length at most \(x/D\), using at most \(1+D\ell_i/x\) cells for a component of length \(\ell_i\). A zero-length component needs one cell. Assign shared endpoints to one cell only, and let \(m_c\) be the cell occupancies. All ordered distinct pairs inside a cell are counted by \(R_D(x)\). Cauchy--Schwarz gives \[
 \frac1D\sum_c m_c(m_c-1)
 \geq \frac D{\#\text{cells}}-1
 \geq\frac1{K/D+\ell/x}-1.
\] This proves (RB.R2). For (RB.R3), the possible origins are covered by the original arcs enlarged backwards by \(L/D\); their total length is at most \(\ell+KL/D\). End of proof.

Consequently, along any sequence with \(K/D\to0\) and \(\ell\to0\), the gap measures converge weakly to \(\delta_0\), every fixed-threshold \(R_D(x)\) diverges, and the local point measure at a uniform origin converges vaguely in probability to the empty measure. The last conclusion follows from (RB.R3) for a compact window containing the support of any test function. It does not rely on an inference from second moments, nor on uniform integrability. The expected number in a window is still exactly \(L\) at every finite dimension.

\subsection{Perturbations much larger than the mean spacing}\label{robustness:perturbations-much-larger-than-the-mean-spacing}

Use the notation of the locality theorem, with odd Pell dimension \(D=(N-1)/2=Kn\), the natural single-kick unitary \(V_N\), its conjugated arithmetic reference \(W_N\), and \[
 \delta_N=\|V_N-W_N\|=O(N^{-\infty}).
\] Let \(\widetilde V_N\) be any other unitary satisfying \[
 \|\widetilde V_N-V_N\|\leq\eta_N,\qquad K\eta_N\to0.       \tag{RB.R4}
\] The same estimate applies to the parity-restricted matrices; \(\widetilde V_N\) need only be defined on that fixed \(D\)-dimensional space. Then all three negative local-statistics conclusions above hold for \(\widetilde V_N\).

\textbf{Proof.} Set \(b_N=\delta_N+\eta_N\). The triangle inequality gives \(\|\widetilde V_N-W_N\|\leq b_N\). For large \(N\), \(b_N<\sin(\pi/K)\). Normal spectral stability and the unitary homotopy argument in the locality theorem put exactly \(n\) eigenvalues in each root cluster of circle-coordinate halfwidth \[
 a_N'\leq\pi^{-1}\arcsin(b_N/2).
\] Thus the spectrum has a cover with total length \[
 \ell_N\leq\frac{2K}{\pi}\arcsin(b_N/2)=O(Kb_N)\to0.       \tag{RB.R5}
\] Apply (RB.R1)--(RB.R3). End of proof.

This statement is uniform over \emph{all} unitary perturbations with the norm bound (RB.R4); their eigenvectors, internal cluster splitting, and probability law are unrestricted. In particular, for any fixed \(b>0\), the sufficient condition \(\eta_N=O(N^{-b})\) holds since \(K=O(\log N)\). When \(0<b<1\), this allowed error is parametrically larger than \(1/D\).

For example, a unitary perturbation with norm \(O(N^{-1/2})\) still obeys \[
 \frac1D\#\{j:S_j>x\}=O_x(K/D+K/\sqrt N),\qquad
 R_D(x)=\Omega_x(\sqrt N/K),                                \tag{RB.R6}
\] where the lower bound is meant for all sufficiently large \(N\). The exact identity \(R_D(x)=n-1\) from the unperturbed locality theorem is not retained under this larger error: individual clusters may have widths larger than a single mean spacing. The covering proof is what replaces that identity.

This gives a precise quantization-stability statement \textbf{only} for a second quantization meeting (RB.R4). The same classical principal canonical map by itself does not imply such an operator-norm bound; leading Fourier-integral amplitudes or subprincipal conventions must not be silently identified.

\subsection{Necessary distortion of any unfolding that restores a nonempty limit}\label{robustness:necessary-distortion-of-any-unfolding-that-restores-a-nonempty-limit}

Let \(F_N\) be an increasing degree-one map of the circle, with Lipschitz constant \(A_N\). The image of a cover of total length \(\ell_N\) is covered by at most \(K\) arcs with total length at most \(A_N\ell_N\). Their positions and occupancies may be arbitrary. Therefore all three collapse conclusions persist whenever \[
 A_N\ell_N\to0.                                            \tag{RB.R7}
\] This is sharper than requiring that every cluster have diameter smaller than \(1/D\) after unfolding.

For a converse necessary bound, suppose the probability that an image window of a fixed unfolded length \(L\) contains a level tends to some \(p_L>0\). By (RB.R3) and \(K/D\to0\), \[
 \liminf A_N\ell_N\geq p_L.                                \tag{RB.R8}
\] This is a necessary condition, not a sufficient recipe for a desired point process. If one uses the original equal-width cover, then \(\ell_N\leq2Ka_N\), giving \[
 \liminf 2KA_Na_N\geq p_L.                                 \tag{RB.R9}
\] For the explicit Gevrey-2 fixed-strength example and fixed \({\kappa_0}\ne0\), there is a constant \(C_{\kappa_0}\) such that \(a_N\leq C_{\kappa_0} N^2e^{-cN^{1/3}}\). Thus any such unfolding must satisfy \[
 A_N\geq c_{L,{\kappa_0}}\frac{e^{cN^{1/3}}}{KN^2}
 \quad\text{for all sufficiently large }N.                 \tag{RB.R10}
\] To see the quantifiers, first make the left probability at least \(p_L/2\) and \(KL/D\leq p_L/4\); then (RB.R3) gives \(2KA_Na_N\geq p_L/4\). Substitute the displayed upper bound on \(a_N\). If a cluster has zero width, no finite Lipschitz unfolding can split its coincident levels; (RB.R8) remains the appropriate general necessary statement. The constant in (RB.R10) is not claimed numerically sharp.
